\documentclass[aspectratio=169,10pt,english]{beamer}

\input{../../../_preambulo_beamer.tex}
\input{../../../_comandos_beamer.tex}

\selectlanguage{english}

\title{MA1001B - Statistical Modeling for Decision Making}
\subtitle{Lesson 03: Central Tendency and Dispersion}
\author{}
\institute{Tecnol\'ogico de Monterrey}
\date{}

\begin{document}

\begin{frame}[plain]
  \vspace{1.2cm}
  {\Large\bfseries MA1001B - Statistical Modeling for Decision Making\par}
  \vspace{0.7cm}
  {\large Lesson 03: Central Tendency and Dispersion\par}
  \vspace{0.7cm}
  {\color{TecRojo}\rule{\textwidth}{1.2pt}}
  \vspace{0.8cm}

  MA1001B Faculty\\[0.3cm]
  Term\\[0.3cm]
  Tecnol\'ogico de Monterrey
\end{frame}

\begin{frame}{The Descriptive Question}
  \begin{alertblock}{Guiding question}
    Which summaries describe a typical processing time and the variability
    among orders?
  \end{alertblock}

  \vspace{0.5em}
  By the end of this lesson, you can:
  \begin{itemize}
    \item calculate and interpret the mean and median;
    \item calculate sample standard deviation and interquartile range;
    \item pair a measure of center with a compatible measure of spread;
    \item explain how one extreme value affects each summary.
  \end{itemize}

  \vfill
  \scriptsize All 240 order records are synthetic. No real company data are used.
\end{frame}

\begin{frame}{Center and Spread}
  \small
  \begin{center}
    \begin{tabular}{p{0.18\textwidth}p{0.31\textwidth}p{0.37\textwidth}}
      \toprule
      \textbf{Summary} & \textbf{Calculation} & \textbf{Question answered} \\
      \midrule
      Mean & $\bar{x}=\frac{1}{n}\sum x_i$ & Where do values balance? \\
      Median & Middle ordered value & Which value divides the ordered data in half? \\
      Sample $s$ & $\sqrt{\frac{\sum(x_i-\bar{x})^2}{n-1}}$ & How far do values vary around the mean? \\
      IQR & $Q_3-Q_1$ & How wide is the middle 50\%? \\
      \bottomrule
    \end{tabular}
  \end{center}

  \vspace{0.8em}
  \begin{exampleblock}{Interpretation rule}
    Report the summaries in context and in the variable's units. Pair mean with
    sample $s$, or median with IQR, so center and spread tell one coherent story.
  \end{exampleblock}
\end{frame}

\begin{frame}[fragile]{Step 1: Mean and Median}
  \begin{columns}[T]
    \begin{column}{0.40\textwidth}
      \small
      \begin{lstlisting}
values = sample[
    "processing_time_minutes"
]

mean = values.mean()
median = values.median()
      \end{lstlisting}

      \begin{block}{Verified output}
        $n=240$ orders\\
        Mean $=45.19$ min\\
        Median $=44.70$ min\\
        Difference $=0.49$ min
      \end{block}
    \end{column}
    \begin{column}{0.60\textwidth}
      \centering
      \includegraphics[width=\textwidth]{../figures/descriptive_statistics_01_center.png}
      \vspace{0.15em}
      \scriptsize Both numbers describe center. Neither number describes spread.
    \end{column}
  \end{columns}
\end{frame}

\begin{frame}{Step 2: Two Measures of Spread}
  \begin{columns}[T]
    \begin{column}{0.39\textwidth}
      \small
      Sample standard deviation:
      \[
        s=\sqrt{\frac{\sum(x_i-\bar{x})^2}{n-1}}=6.38\text{ min}
      \]

      Quartile interval:
      \[
        Q_1=40.53,\quad Q_3=49.01
      \]
      \[
        \mathrm{IQR}=Q_3-Q_1=8.48\text{ min}
      \]

      $s$ measures variation around the mean. The IQR spans the middle half of
      the ordered observations.
    \end{column}
    \begin{column}{0.61\textwidth}
      \centering
      \includegraphics[width=\textwidth]{../figures/descriptive_statistics_02_spread.png}
    \end{column}
  \end{columns}
\end{frame}

\begin{frame}{Step 3: Sensitivity to One Extreme Value}
  \small
  \textbf{Order 300193 changes from 65.22 to 180.00 minutes.}

  \vspace{0.1em}
  \centering
  \includegraphics[width=0.80\textwidth]{../figures/descriptive_statistics_03_sensitivity.png}

  \vspace{0.1em}
  \begin{tabular}{lrrrr}
    \toprule
    \textbf{Summary} & Mean & Median & Sample $s$ & IQR \\
    \textbf{Change (min)} & $+0.48$ & $+0.00$ & $+4.34$ & $+0.00$ \\
    \bottomrule
  \end{tabular}

  \vspace{0.1em}
  \scriptsize Displayed changes are rounded to two decimals. Sensitivity alone
  does not diagnose an anomaly.
\end{frame}

\begin{frame}{Choosing Descriptive Summaries}
  \small
  \textbf{Mean with sample standard deviation}\par\vspace{0.2em}
  This pair uses every observation. It often describes a roughly symmetric
  distribution well when no extreme values dominate the summaries.

  \vspace{0.7em}
  \textbf{Median with interquartile range}\par\vspace{0.2em}
  This pair depends on ordered positions. It usually resists the influence of a
  small number of extreme values better.

  \vspace{0.8em}
  \begin{alertblock}{A defensible description}
    Inspect the distribution, state the units, and explain the choice. These
    pairings are useful guidelines rather than automatic rules.
  \end{alertblock}

  A numeric identifier such as \texttt{order\_id} still labels a record. It is
  not a quantitative measurement to summarize with a mean or standard deviation.
\end{frame}

\begin{frame}{Concept Check}
  \begin{enumerate}
    \item What different questions do the mean and sample standard deviation
      answer about processing time?
    \item Using $Q_1=40.53$ and $Q_3=49.01$, calculate and interpret the IQR.
    \item Which center-spread pair better describes data affected by a small
      number of extreme delays, and why?
    \item Why does the 180-minute value require more evidence before anyone
      labels it an anomaly or recording error?
  \end{enumerate}

  \vspace{0.5em}
  \begin{block}{Connection to Lesson 04}
    Boxplots and explicit detection rules add a reproducible way to flag values
    for investigation. Business context still determines what the flag means.
  \end{block}

  \vfill
  \scriptsize Source: OpenStax, \textit{Introductory Business Statistics 2e},
  Chapter 2, Sections 2.2, 2.3, and 2.7. CC BY-NC-SA 4.0.
\end{frame}

\appendix



\end{document}
